% ==========================================================================
% Introduction
% ==========================================================================

\section{Introduction}

Modern oncology increasingly relies on computational models for
predicting tumor evolution and optimizing anticancer therapies. In
particular, recent advances in Large Language Models (LLMs) and
Retrieval-Augmented Generation (RAG) systems have opened the possibility
of developing interactive frameworks for Language--Driven Therapy Design
(LDTD), where clinicians can explore treatment scenarios and control
predictive simulations through natural language \cite{reka2026conversational}. However, practical
implementation of such systems requires simulation engines operating in
near real-time, which remains infeasible for complex mechanistic tumor
models involving high-dimensional spatial dynamics and computationally
expensive partial differential equations.

A natural solution is the use of surrogate models that approximate the
behavior of detailed mechanistic simulations at a much lower
computational cost. Nevertheless, aggressive model reduction -
especially the elimination of spatial information - introduces
substantial challenges. Reduced ordinary differential equation (ODE)
models frequently fail to capture important effects arising from tumor
geometry, spatial heterogeneity, and mismatches between tumor and
treatment-field geometry. As a consequence, unresolved spatial dynamics
are often absorbed into distorted effective model parameters, leading to
biased predictions and reduced interpretability.

In this work, we investigate whether simple physics-informed surrogate
models are capable of partially reconstructing the effects of neglected
spatial dynamics despite operating in a drastically simplified
non-spatial representation. We focus on radiotherapy-driven tumor
dynamics (e.g. \cite{Balcerak2025GliODIL}, \cite{sun2025physicsinformed}, \cite{Sujitha2023RadiotherapyBrainTumor}) in scenarios where the irradiation field only partially
overlaps with the tumor volume or becomes spatially shifted with respect
to the evolving tumor geometry during repeated treatment fractions. Such
situations naturally generate unresolved effects including asymmetric
survival regions, delayed recolonization of irradiated areas, outward
invasive expansion, and spatially heterogeneous regrowth dynamics.
The presented framework is intended as a controlled computational
testbed for investigating whether physics-informed surrogate models can
reconstruct unresolved spatial effects eliminated during PDE-to-ODE
reduction. From the Scientific Machine Learning perspective, the
PI-NODE architecture may also be interpreted as a reduced-order
UDE-based \cite{rackauckas2019diffeqflux} closure model designed to compensate for unresolved spatial
dynamics introduced by aggressive PDE-to-ODE reduction including spatially heterogeneous radiation damage accumulation,
asymmetric post-treatment survival patterns, and delayed
recolonization fronts.

The experiments presented should therefore be interpreted primarily as a
methodological investigation of closure-aware surrogate modeling for
future interactive therapy-planning systems and Language–Driven Therapy
Design frameworks requiring computationally efficient predictive models.
